\scanpage{535}%
\mspage{446}%
% archive sheet number: 534
cardinal $\leq\pi_0$, et $\deleted{\unsure{I|J}}J_{\deleted{\unsure{\geq}}{\geq j_0}}$ est grd devant $\pi_0$, on voit qu'il existe $J'\subset{}_{j_0}\backslash J_{\deleted{\unsure{\geq j_0}}}$, cofinal ds $\mathord{}_{j_0}\backslash J_{\deleted{\unsure{\geq j_0}}}$\inserted{donc ds $J$}, tel que pour $j\in J'$, les $\tau_{j,j_0}^\natural$ soient égaux à un $\deleted{\unsure{=|\tau}}\ \tau':x_0^\natural\xrightarrow{\sim}x_0^\natural$.

Donc si $\deleted{j'\leq j''}\ j',j''\in J'$, avec $j'\leq j''$, on a

\[\underbrace{\tau_{j''j_0}}_{\tau}=\tau_{j''j'}\underbrace{\tau_{j',j}}_{\tau},\]

\[\text{d'où}\qquad\tau_{j''j'}=\mathrm{id}_{x_0^\natural},\]

puisque la flèche $\tau$ \deleted{\unsure{sont des}} est un iso. Dans \deleted{\unsure{si|pour}} le système inductif réindexé par $J'$, \deleted{\unsure{on a}} la condition $4^\circ)$ est \uline{remplacée par la con}dition plus stricte

\[\boxed{4^\circ\text{bis})\quad\deleted{\text{\unsure{pour|que}}}\quad\tau_{j',j}^\natural=\mathrm{id}_{x_0^\natural}.}\]

Mais la condition $\deleted{3}2^\circ)$ devient alors, grâce à 4.15.7.4\inserted{\unsure{et compte tenu de|et condition de} $3^\circ,4$},

\[\boxed{\deleted{3}2^\circ\text{bis})\quad\forall j,\ \text{existe }j'\geq j\ \text{avec }\tau_{j',j}=p}. \]

\deleted{Donc}, OPS $1^\circ,2^\circ\mathrm{bis},3^\circ,4^\circ\mathrm{bis}$. (on désigne à nouveau par $J$, \inserted{\unsure{ledit|ci-dt} $J'$,} l'ens. ordonné grd devant $\pi_0$ indexant.).

\deleted{\unsure{Distinguons deux cas :} \deleted{\unsure{Soit $K$ l'ens ord.}} $a)$ \deleted{$J$} $I$ admet un objet final -- alors $I$ aussi, \unsure{Notons que comme}}

Soit $K$ l'ens ordonné défini par
